\subsection{Quantum search\label{SEARCH}}

In our security proof for quantum money, an important step will be to
\emph{amplify} a counterfeiter\ who copies a banknote $\$$ with any
non-negligible fidelity to a counterfeiter who copies $\$$ almost perfectly.
\ Taking the contrapositive, this will imply that to rule out the former sort
of counterfeiter, it suffices to rule out the latter.

In this section, we first review two variants of Grover's search algorithm~\cite{grover}\ that are useful for amplifying the fidelity of quantum states.
\ We then introduce a variant that combines the advantages of both.

Assume we are given a pure initial state $\left\vert \operatorname*{Init}%
\right\rangle $, in some Hilbert space $\mathcal{H}$. Our goal is to map
$\left\vert \operatorname*{Init}\right\rangle $ to a final state $\left\vert
\Psi\right\rangle $ that lies in (or close to) a \textquotedblleft good
subspace\textquotedblright\ $G\leq\mathcal{H}$. We have oracle access to two
unitary transformations:

\begin{itemize}
\item $U_{\operatorname*{Init}}$, which maps $\left\vert \operatorname*{Init}%
\right\rangle $ to $-\left\vert \operatorname*{Init}\right\rangle $, and acts
as the identity on all $\left\vert v\right\rangle $ orthogonal to $\left\vert
\operatorname*{Init}\right\rangle $.

\item $U_{G}$, which maps $\left\vert v\right\rangle $ to $-\left\vert
v\right\rangle $ for all $\left\vert v\right\rangle \in G$, and acts as the
identity on all $\left\vert v\right\rangle $ orthogonal to $G$.
\end{itemize}

We are promised that the fidelity of the initial state with $G$,%
\[
F\left(  \left\vert \operatorname*{Init}\right\rangle ,G\right)
=\max_{\left\vert \psi\right\rangle \in G}\left\langle \operatorname*{Init}%
|\psi\right\rangle ,
\]
is at least some $\varepsilon>0$.

In this scenario, \emph{provided }$F\left(  \left\vert \operatorname*{Init}%
\right\rangle ,G\right)  $\emph{\ is known}, the amplitude amplification
framework of Brassard, H\o yer, Mosca, and Tapp~\cite{bhmt} lets us prepare a
state close to $G$ using only $\Theta(1/\varepsilon)$ iterations:

\begin{lemma}
[Amplitude Amplification~\cite{bhmt}]\label{aa}Write $\left\vert
\operatorname*{Init}\right\rangle $ as $\sin\theta\left\vert
\operatorname*{Good}\right\rangle +\cos\theta\left\vert \operatorname*{Bad}%
\right\rangle $, where $\left\vert \operatorname*{Good}\right\rangle $ is the
unit vector formed by projecting $\left\vert \operatorname*{Init}\right\rangle
$ onto $G$, and $\left\vert \operatorname*{Bad}\right\rangle $ is orthogonal
to $\left\vert \operatorname*{Good}\right\rangle $. Then by using $O\left(
T\right)  $ oracle calls\ to $U_{\operatorname*{Init}}$ and $U_{G}$, we can
prepare the state%
\[
\left\vert \Phi_{T}\right\rangle :=\sin\left[  \left(  2T+1\right)
\theta\right]  \left\vert \operatorname*{Good}\right\rangle +\cos\left[
\left(  2T+1\right)  \theta\right]  \left\vert \operatorname*{Bad}%
\right\rangle\,.
\]

\end{lemma}

Note that Grover's algorithm is simply a special case of \expref{Lemma}{aa}, where
$\left\vert \operatorname*{Init}\right\rangle $ is the uniform superposition
over $N$ basis states $\left\vert 1\right\rangle ,\ldots,\left\vert
N\right\rangle $, and $G$ is the subspace spanned by \textquotedblleft
marked\textquotedblright\ states.

However, \expref{Lemma}{aa}\ has an annoying drawback, which it shares with
ordinary Grover search. Namely, the algorithm does \emph{not} converge
monotonically toward the target subspace $G$, but could instead
\textquotedblleft wildly overshoot it,\textquotedblright\ cycling around the
$2$-dimensional\ subspace spanned by $\left\vert \operatorname*{Bad}%
\right\rangle $ and $\left\vert \operatorname*{Good}\right\rangle $. If we
know the fidelity $F(\vert \operatorname*{Init}\rangle
,G)$ in advance (rather than just a lower bound\ on the fidelity),
\emph{or} if we can prepare new copies of $\left\vert \operatorname*{Init}%
\right\rangle $ \textquotedblleft free of charge\textquotedblright\ in case of
failure, then this overshooting is not a serious problem. Alas, neither of
those conditions will hold in our application.

Fortunately, for independent reasons, in 2005 Tulsi, Grover, and Patel~\cite{tgp} introduced a new quantum search algorithm that \emph{does}
guarantee monotonic convergence toward $G$, by alternating unitary
transformations with measurements. (Their algorithm was later simplified and
improved by Chakraborty, Radhakrishnan, and Raghunathan~\cite{crr}.)

\begin{lemma}
[Fixed-Point Quantum Search~\cite{tgp,crr}]\label{fixedpoint}By using $T$
oracle calls\ to $U_{\operatorname*{Init}}$ and $U_{G}$, we can prepare a
state $\left\vert \Psi\right\rangle $ such that $F\left(  \left\vert
\Psi\right\rangle ,G\right)  \geq1-\exp(  -T\varepsilon^{2})$.
\end{lemma}

Rearranging, \expref{Lemma}{fixedpoint}\ lets us prepare a state $\left\vert
\Psi\right\rangle $ such that $F(\vert \Psi\rangle
,G)  \geq1-\delta$ using 
\[
T=O\left(\frac{1}{\varepsilon^{2}}\log
\frac{1}{\delta}\right)
\]
iterations. On the positive side, the dependence
on $1/\delta$ in this bound is logarithmic: we get not only monotonic
convergence toward $G$, but \emph{exponentially-fast} convergence. On the
negative side, notice that the dependence on $\varepsilon$ has worsened from
$1/\varepsilon$ to $1/\varepsilon^{2}$---negating the quadratic speedup that
was the original point of quantum search!

In the rest of this section, we give a \textquotedblleft
hybrid\textquotedblright\ quantum search algorithm that combines the
advantages of Lemmas \ref{aa}\ and \ref{fixedpoint}---i.e., it converges
monotonically toward the target subspace $G$ (rather than \textquotedblleft
overshooting\textquotedblright\ $G$), but \emph{also} achieves a quadratic
speedup. In the context of our security proof for quantum money, this hybrid
algorithm will lead to a quadratically-better (and in fact, tight) lower bound
on the number of queries that a counterfeiter needs to make, compared to what
we would get from using \expref{Lemma}{fixedpoint} alone. While this quadratic
improvement is perhaps only of moderate interest, we include the algorithm in
the hope that it will find other applications.

We first give a technical lemma needed to analyze our algorithm.

\begin{lemma}
\label{interval}For all $L,\beta,\eta,\gamma$, there are at most $\left(
L/\beta+1\right)  \left(  2\eta+1\right)  $ integers $T\in\left\{
0,\ldots,L\right\}$ such that $\left\vert T-\left(  \beta n+\gamma\right)
\right\vert <\eta$ for some integer $n$.
\end{lemma}

\begin{proof}
The real interval $\left[  0,L\right]  $ can intersect at most $L/\beta
+1$ intervals $\left(  \beta n+\gamma-\eta,\beta n+\gamma+\eta\right)  $, and
each such interval can contain at most $2\eta+1$ integer points.
\end{proof}

We now give our hybrid of Lemmas \ref{aa}\ and \ref{fixedpoint}.

\begin{theorem}
[Faster Fixed-Point Search]\label{combined}Let $\delta\geq2\varepsilon$.
\ Then by using \[
O\left(  \frac{\log1/\delta}{\varepsilon\delta^{2}}\right)
\]
oracle calls\ to $U_{\operatorname*{Init}}$ and $U_{G}$, we can prepare a
state $\rho$ such that $F\left(  \rho,G\right)  \geq1-\delta$.
\end{theorem}

\begin{proof}
Let $\xi:=\arcsin\varepsilon$; note that $\varepsilon\leq\xi\leq\frac{\pi}%
{2}\varepsilon$. Also let $L:=\left\lceil 100/\xi\right\rceil $ and
$R:=({25}/{\delta^{2}})\left(  2+\log({1}/{\delta})\right)  $. Then the
algorithm is as follows:

\begin{enumerate}
\item[(1)] Choose an integer $T\in\left\{  0,\ldots,L\right\}  $ uniformly at random.

\item[(2)] Apply $T$ iterations of amplitude amplification with $\left\vert
\operatorname*{Init}\right\rangle $ as the initial state and $G$ as the
target subspace (as in \expref{Lemma}{aa}), to obtain a state $\left\vert \Phi
_{T}\right\rangle $.

\item[(3)] Apply $R$ iterations of fixed-point quantum search with
$\left\vert \Phi_{T}\right\rangle $ as the initial state and $G$ as the
target subspace (as in \expref{Lemma}{fixedpoint}), to obtain a state $\left\vert
\Psi_{T}\right\rangle $.
\end{enumerate}

The final output of the above algorithm is
\[
\rho=\operatorname*{E}_{T\in\left\{  0,\ldots,L\right\}  }\bigl[  \left\vert
\Psi_{T}\right\rangle \left\langle \Psi_{T}\right\vert \bigr]  .
\]
Also, the total number of oracle calls to $U_{\operatorname*{Init}}$ and
$U_{G}$ is%
\[
O\left(  TR\right)  =O\left(  \frac{\log1/\delta}{\varepsilon\delta^{2}%
}\right)  .
\]
(The reason this number scales like $TR$ rather than $T+R$ is that, in step
(3), each time we reflect about the initial state $\left\vert \Phi
_{T}\right\rangle $ we need to rerun step (2). Thus, we need $\Theta\left(
T\right)  $ oracle calls within each of the $R$ iterations.)

By \expref{Lemma}{aa}, after step (2) we have a state $\left\vert \Phi
_{T}\right\rangle $ such that
\[
F\bigl(  \left\vert \Phi_{T}\right\rangle ,G\bigr)  =\bigl\vert \left\langle
\Phi_{T}|\operatorname*{Good}\right\rangle \bigr\vert =\bigl\vert \sin\left[
\left(  2T+1\right)  \xi\right]  \bigr\vert .
\]
So for any $\alpha\in\left(  0,1\right)  $,%
\begin{align*}
\Pr_{T\in\left\{  0,\ldots,L\right\}  }\bigl[  F\left(  \left\vert \Phi
_{T}\right\rangle ,G\right)  <\alpha\bigr]   &  =\Pr_{T\in\left\{
0,\ldots,L\right\}  }\bigl[  \left\vert \sin\left[  \left(  2T+1\right)
\xi\right]  \right\vert <\alpha\bigr]  \\
&  =\Pr_{T\in\left\{  0,\ldots,L\right\}  }\bigl[  \exists n\in\mathbb{Z}%
:\left\vert \left(  2T+1\right)  \xi-\pi n\right\vert <\arcsin\alpha\bigr]
\\
&  \leq\frac{\left(  \frac{L}{\pi/2\xi}+1\right)  \left(  \frac{\arcsin\alpha
}{\xi}+1\right)  }{L+1}\\
&  \leq\frac{2}{\pi}\arcsin\alpha+\frac{2\xi}{\pi}+\frac{\arcsin\alpha}%
{100}+\frac{\xi}{100}\\[1ex]
&  \leq1.02\left(  \alpha+\varepsilon\right)  ,
\end{align*}
where the third line uses \expref{Lemma}{interval}.

Now assume $F(\vert \Phi_{T}\rangle ,G)  \geq\alpha$.
\ Then by \expref{Lemma}{fixedpoint}, after step (3) we have a state $\left\vert
\Psi_{T}\right\rangle $ such that%
\[
F(\vert \Phi_{T}\rangle ,G)  \geq1-\exp\left(
-R\alpha^{2}\right)  .
\]
Let us now make the choice $\alpha:=\delta/5$. Then by the union bound, the
\textquotedblleft average\textquotedblright\ output $\rho=\operatorname*{E}%
_{T}\left[  \left\vert \Psi_{T}\right\rangle \left\langle \Psi_{T}\right\vert
\right]  $ satisfies%
\begin{align*}
1-F\left(  \rho,G\right)   &  \leq1.02\left(  \alpha+\varepsilon\right)
+\exp\left(  -R\alpha^{2}\right) \\
&  \leq0.204\delta+0.51\delta+\exp\left(  -\frac{R\delta^{2}}{25}\right) \\
&  <\delta\,.\qedhere
\end{align*}
\end{proof}
